% Canonical reconciliation of the independently delivered correlation reports.
\providecommand{\stconvT}{\mathbb T}
\providecommand{\stconvId}{\mathsf I}
\providecommand{\stconvG}{\mathcal G}
\providecommand{\stconvBell}{\mathcal B}
\providecommand{\stconvCT}{\operatorname{CT}}
\providecommand{\stconvFp}{\operatorname{Fp}}
\providecommand{\FP}{\operatorname{FP}}
The single-function jet calculus extends to products once the endpoint
normalization is fixed. The four independent continuations
\cite{stconv:algebra-report,stconv:shifted-report,stconv:closure-report,stconv:calculus-report}
agree on the scale-one finite parts used below. Their common proof is
organized through Fourier normalization, singular Laurent coefficients,
convolution closure, translated products and normalized primitive lifting.
The contact terms belong to the differentiated extension in the next chapter.

\section{Fourier conventions and the Hurwitz semigroup}
\label{stconv:sec:semigroup}
We work on $\stconvT=\mathbb R/\mathbb Z$, with normalized Lebesgue measure.
For a periodic distribution $T$, use the convention
\[
 \widehat T(k)=\langle T,e^{-2\pi i kx}\rangle,\qquad k\in\Z.
\]
The constant function $1$ and the point mass $\delta_0$ both have mean
one, but their nonzero Fourier coefficients are respectively $0$ and $1$.
Thus
\[
 \stconvId=\delta_0-1
\]
is the convolution identity on the space of mean-zero distributions. In
particular, $\stconvId*T=T$ whenever $\widehat T(0)=0$.

The logarithm used throughout is
\begin{equation}\label{stconv:eq:logbranch}
 \lambda_k=\log(2\pi i k)
 =\log(2\pi|k|)+\frac{\pi i}{2}\operatorname{sgn}k,
 \qquad k\ne0.
\end{equation}
All powers below refer to this fixed branch.

\begin{definition}\label{stconv:def:R}
For $\alpha\in\C$, define $R_\alpha$ by
\begin{equation}\label{stconv:eq:Rfourier}
 \widehat R_\alpha(0)=0,\qquad
 \widehat R_\alpha(k)=e^{-\alpha\lambda_k}\quad(k\ne0).
\end{equation}
Equivalently, its Fourier series is
$\sum_{k\ne0}(2\pi i k)^{-\alpha}e^{2\pi i kx}$ in the distribution sense.
\end{definition}

\begin{lemma}[Entire distribution family]\label{stconv:lem:entire}
The family $\alpha\mapsto R_\alpha$ is entire as a periodic-distribution
family. It is represented by an absolutely convergent Fourier series for
$\Rea\alpha>1$. For $\Rea\alpha>0$, it is the locally integrable periodic
extension of
\begin{equation}\label{stconv:eq:Rhurwitz}
 R_\alpha(x)=\frac{\zeta(1-\alpha,x)}{\Gamma(\alpha)},\qquad0<x<1.
\end{equation}
\end{lemma}
\begin{proof}
On each compact set of parameters, the coefficients in
\eqref{stconv:eq:Rfourier}, together with any fixed number of parameter derivatives,
have at most polynomial growth in $|k|$. Fourier coefficients of a smooth
test function decay faster than every polynomial. Hence the paired Fourier
series and all its parameter derivatives converge locally uniformly. This
proves the first assertion; the absolute-convergence assertion follows
immediately from $\sum k^{-\Rea\alpha}$.

For $\Rea\alpha>1$, equation~\eqref{stconv:eq:Rhurwitz} is precisely Hurwitz's
Fourier formula. To continue to $\Rea\alpha>0$, use
\[
 \zeta(1-\alpha,x)=x^{\alpha-1}+\zeta(1-\alpha,1+x).
\]
On compact parameter subsets of that half-plane, the first term and its
parameter derivatives are dominated by integrable powers of $x$ times
powers of $|\log x|$. The second term is regular in $x\in[0,1]$; any
apparent parameter singularity at $\alpha=0$ is outside the stated open
half-plane. The two holomorphic distribution families agree in an open
subregion, and hence agree throughout the half-plane.
\end{proof}

\begin{theorem}[Semigroup and differential ladder]\label{stconv:thm:semigroup}
For all $\alpha,\beta\in\C$,
\begin{equation}\label{stconv:eq:semigroup}
 R_\alpha*R_\beta=R_{\alpha+\beta},\qquad
 DR_\alpha=R_{\alpha-1},\qquad R_0=\stconvId.
\end{equation}
If $J_{\alpha,r}=\partial_\alpha^rR_\alpha$, then
\begin{equation}\label{stconv:eq:mixedjets}
 J_{\alpha,r}*J_{\beta,s}=J_{\alpha+\beta,r+s},\qquad
 D^kJ_{\alpha,r}=J_{\alpha-k,r}.
\end{equation}
\end{theorem}
\begin{proof}
Periodic distributions can be convolved: equivalently, their
polynomially growing Fourier coefficients can be multiplied. For $k\ne0$,
\[
 e^{-\alpha\lambda_k}e^{-\beta\lambda_k}
 =e^{-(\alpha+\beta)\lambda_k},\qquad
 (2\pi i k)e^{-\alpha\lambda_k}=e^{-(\alpha-1)\lambda_k}.
\]
Both sides have zero zeroth coefficient. Fourier uniqueness proves
\eqref{stconv:eq:semigroup}. Differentiating the locally uniformly paired
Fourier series gives~\eqref{stconv:eq:mixedjets}.
\end{proof}

For real $\alpha,\beta>1$, the first identity, expressed using
\eqref{stconv:eq:Rhurwitz}, is the nonsingular law in Sun~\cite[Remark 3.2]{stconv:Sun}.
The new use here is at its singular integer parameters and their spectral
jets. At positive integers,
\[
 R_m(x)=-\frac{B_m(\{x\})}{m!}\quad(m\ge1),
\]
up to immaterial point values of the sawtooth for $m=1$. At negative
integers,
\begin{equation}\label{stconv:eq:negativeR}
 R_{-k}=\delta_0^{(k)}\quad(k\ge1).
\end{equation}
The quotient in~\eqref{stconv:eq:Rhurwitz} vanishes pointwise at $\alpha=-k$
away from the integers, while the distribution is nonzero. This is the
basic reason contact terms must be retained.

\section{Canonical finite parts of the Stieltjes functions}
\label{stconv:sec:finiteparts}
The generalized Stieltjes functions are defined for $x>0$ by
\begin{equation}\label{stconv:eq:laurent}
 \zeta(1-z,x)=-\frac1z+\sum_{n=0}^\infty\gamma_n(x)\frac{z^n}{n!}.
\end{equation}
Thus $\gamma_0(x)=-\psi(x)$. The Hurwitz shift relation gives
\begin{equation}\label{stconv:eq:stshift}
 \gamma_n(x)=\frac{\log^n x}{x}+\gamma_n(1+x).
\end{equation}
Consequently, periodic extension of $\gamma_n$ is not locally integrable
at $0$; a convention is necessary.

\begin{definition}[Canonical periodic finite part]\label{stconv:def:G}
For a smooth periodic test function $\phi$, define
\begin{align}
 \langle G_n,\phi\rangle
 &=\int_0^1\left(\gamma_n(x)\phi(x)
       -\frac{\log^n x}{x}\phi(0)\right)\dd x\label{stconv:eq:Gdefinition}\\
 &=\lim_{\epsilon\downarrow0}\left\{
       \int_\epsilon^1\gamma_n(x)\phi(x)\dd x
       +\frac{\log^{n+1}\epsilon}{n+1}\phi(0)\right\}.
       \label{stconv:eq:Gcutoff}
\end{align}
\end{definition}
The integrand in~\eqref{stconv:eq:Gdefinition} is integrable: its singular part is
$\log^n x\,[\phi(x)-\phi(0)]/x$, which is $O(|\log x|^n)$.
No arbitrary finite multiple of $\delta_0$ is added.

\begin{lemma}[The distributional Laurent expansion]\label{stconv:lem:Laurentdist}
Let
\[
 Z(s)=\Gamma(1-s)R_{1-s}.
\]
It is a meromorphic periodic-distribution family, agreeing with the
periodic locally integrable Hurwitz zeta for $\Rea s<1$. Near $z=0$,
\begin{equation}\label{stconv:eq:ZLaurent}
 Z(1-z)=\frac{\stconvId}{z}+\sum_{n=0}^\infty G_n\frac{z^n}{n!}.
\end{equation}
Every $G_n$ has mean zero.
\end{lemma}
\begin{proof}
The meromorphic continuation follows from Lemma~\ref{stconv:lem:entire}. For
$\Rea z>0$, split the pairing using
$\zeta(1-z,x)=x^{z-1}+\zeta(1-z,1+x)$. The first term has expansion
\[
 \int_0^1x^{z-1}\phi(x)\dd x
 =\frac{\phi(0)}z+
   \sum_{n\ge0}\frac{z^n}{n!}
   \int_0^1\frac{\log^n x}{x}[\phi(x)-\phi(0)]\dd x.
\]
The second term contributes $-z^{-1}\int_0^1\phi$ and the coefficients
$\int_0^1\gamma_n(1+x)\phi(x)\dd x$. These expansions converge locally
near zero after pairing, by the same endpoint domination used above.
Combining them proves~\eqref{stconv:eq:ZLaurent} with the exact subtraction in
Definition~\ref{stconv:def:G}. Since $\widehat Z(s)(0)=0$, every coefficient has
mean zero. Equivalently, direct integration over $[1,2]$ gives
$\int_1^2\zeta(1-z,x)\dd x=-1/z$, whose regular coefficients vanish.
\end{proof}

\begin{remark}[Two different residues]
The scalar function in~\eqref{stconv:eq:laurent} has residue $-1$ as a function
of $z$, whereas the periodic-distribution family has residue
$\delta_0-1$. Analytic continuation across a nonintegrable endpoint adds
a point-supported residue. Treating the scalar residue as the whole
periodic residue gives wrong convolution constants.
\end{remark}

\section{All-index convolution closure and exact normal forms}
\label{stconv:sec:algebra}
Write $\stconvG(z)=\sum_{n\ge0}G_nz^n/n!$, holomorphic as a distribution family
for $|z|<1$. By~\eqref{stconv:eq:ZLaurent},
\begin{equation}\label{stconv:eq:masterG}
 \stconvId+z\stconvG(z)=\Gamma(1+z)R_z.
\end{equation}
Let $X=G_0$. For a complete exponential Bell polynomial we use
\begin{equation}\label{stconv:eq:Belldef}
 \exp\left(\sum_{j\ge1}b_j\frac{z^j}{j!}\right)
 =\sum_{r\ge0}\stconvBell_r(b_1,\ldots,b_r)\frac{z^r}{r!}.
\end{equation}
When its arguments are distributions, all products in this definition
are convolution products and its scalar unit is $\stconvId$.

\begin{theorem}[One-generator convolution algebra]\label{stconv:thm:Bell}
For every $n\ge0$,
\begin{equation}\label{stconv:eq:BellG}
 (n+1)G_n=\stconvBell_{n+1}(b_1,\ldots,b_{n+1}),\qquad
 b_1=X,\quad b_j=(-1)^j(j-1)!\zeta(j)\stconvId\ (j\ge2).
\end{equation}
The span of $\stconvId,G_0,G_1,\ldots$ is a convolution algebra isomorphic to
the polynomial algebra $\C[X]$. These distributions form a basis of that
algebra, with the unit listed first.
\end{theorem}
\begin{proof}
For $k\ne0$, the Fourier transform of~\eqref{stconv:eq:masterG} is
\[
 1+z\widehat\stconvG(z)(k)
 =\Gamma(1+z)e^{-z\lambda_k}
 =\exp\left[(-\gamma-\lambda_k)z+
       \sum_{j\ge2}\frac{(-1)^j\zeta(j)}jz^j\right].
\]
The linear coefficient is
$\widehat X(k)=-\gamma-\lambda_k$. Comparing with
\eqref{stconv:eq:Belldef} proves~\eqref{stconv:eq:BellG}; the zero Fourier coefficient
of each distributional expression vanishes. The leading term in $G_n$
is $X^{*(n+1)}/(n+1)$, so the listed distributions span exactly the
polynomial algebra generated by $X$.

There are no polynomial relations among its powers. Indeed, if
$P(X)=0$ distributionally, then
$P(-\gamma-\log(2\pi k)-i\pi/2)=0$ for every positive integer $k$.
These are infinitely many distinct complex numbers, and hence $P$ is the
zero polynomial. This also proves the basis assertion.
\end{proof}

This algebraic freeness concerns distributions as functions of the
variable $x$. It is not an independence theorem about their values at a
fixed rational point, nor about zeta or Stieltjes constants as numbers.

The first Bell forms are
\begin{align}
 2G_1&=X^{*2}+\zeta(2)\stconvId,\label{stconv:eq:Belllow}\\
 3G_2&=X^{*3}+3\zeta(2)X-2\zeta(3)\stconvId,\notag\\
 4G_3&=X^{*4}+6\zeta(2)X^{*2}-8\zeta(3)X
                +[3\zeta(2)^2+6\zeta(4)]\stconvId.\notag
\end{align}
They already imply many exact integrals. An efficient all-index
multiplication rule is supplied by the following Gamma quotient.

\begin{theorem}[Bivariate multiplication law]\label{stconv:thm:product}
Define, near $(u,v)=(0,0)$,
\begin{align}
 B(u,v)&=\frac{\Gamma(1+u)\Gamma(1+v)}{\Gamma(1+u+v)}\label{stconv:eq:Buv}\\
 &=\exp\left(\sum_{j\ge2}\frac{(-1)^j\zeta(j)}j
          [u^j+v^j-(u+v)^j]\right).\notag
\end{align}
Then
\begin{align}
 uv\,\stconvG(u)*\stconvG(v)={}&[B(u,v)-1]\stconvId\label{stconv:eq:Gproduct}\\
 &+B(u,v)(u+v)\stconvG(u+v)-u\stconvG(u)-v\stconvG(v).\notag
\end{align}
The right side is divisible by $uv$ as a convergent distribution-valued
power series. In particular,
\begin{equation}\label{stconv:eq:topindex}
 G_m*G_n-\frac{m+n+2}{(m+1)(n+1)}G_{m+n+1}
       \in\operatorname{span}_{\C}(\{\stconvId\}\cup\{G_j:0\le j<m+n\}),
\end{equation}
where the lower Stieltjes span is empty when $m+n=0$. The index
$m+n$ is absent because the Gamma quotient has no linear term. Every
coefficient is a polynomial over $\Q$ in zeta values at positive
integers. Its complete value is obtained by multiplying the coefficient
of $u^{m+1}v^{n+1}$ in~\eqref{stconv:eq:Gproduct} by $m!n!$.
\end{theorem}
\begin{proof}
Convolve~\eqref{stconv:eq:masterG} at parameters $u$ and $v$ and apply
Theorem~\ref{stconv:thm:semigroup}. The result is
\[
 [\stconvId+u\stconvG(u)]*[\stconvId+v\stconvG(v)]
 =B(u,v)[\stconvId+(u+v)\stconvG(u+v)].
\]
Subtract the unit and the two linear summands. Both restrictions $u=0$
and $v=0$ vanish identically, proving divisibility. Coefficient extraction
is legitimate in a sufficiently small bidisc. The largest-index term
comes from $(u+v)\stconvG(u+v)$ with the constant coefficient of $B$; its
coefficient is
$\frac{m!n!}{(m+n+1)!}\binom{m+n+2}{m+1}$, which equals the number in
\eqref{stconv:eq:topindex}. The exponential expansion of $B$ proves the
zeta-polynomial assertion.
\end{proof}

Four useful multiplication rules are
\begin{align}
 G_0*G_0&=2G_1-\zeta(2)\stconvId,\label{stconv:eq:product00}\\
 G_0*G_1&=\frac32G_2-\zeta(2)G_0+\zeta(3)\stconvId,\label{stconv:eq:product01}\\
 G_0*G_2&=\frac43G_3-2\zeta(2)G_1+2\zeta(3)G_0-2\zeta(4)\stconvId,\label{stconv:eq:product02}\\
 G_1*G_1&=G_3-2\zeta(2)G_1+2\zeta(3)G_0-\frac14\zeta(4)\stconvId.
 \label{stconv:eq:product11}
\end{align}
The last simplification uses $\zeta(2)^2=5\zeta(4)/2$. The verification
catalog retains independent formal symbols for zeta values until such a
specific relation is explicitly applied.

\section{Ordinary integral realizations and explicit Stieltjes identities}
\label{stconv:sec:integrals}
The singular support of $G_n$ is contained in $\{0\}$. Thus its
convolution with another $G_m$ is smooth away from $0$: in a local
convolution integral at $0<x<1$, the two singular factors occur at the
distinct integration points $t=0$ and $t=x$. We now give a direct formula
that also proves this local realization without requiring a product of
distributions at the same point.

\begin{theorem}[Absolutely convergent subtracted integrals]\label{stconv:thm:integral}
Put $a_n(t)=\log^n(t)/t$. For $m,n\ge0$ and $0<x<1$, define
\begin{align}
 C_{m,n}(x)={}&\int_0^x\big[\gamma_m(t)\gamma_n(x-t)
          -\gamma_n(x)a_m(t)-\gamma_m(x)a_n(x-t)\big]\dd t
          \label{stconv:eq:Cmn}\\
 &+\int_x^1\gamma_m(t)\gamma_n(1+x-t)\dd t\notag\\
 &+\gamma_n(x)\frac{\log^{m+1}x}{m+1}
  +\gamma_m(x)\frac{\log^{n+1}x}{n+1}.\notag
\end{align}
Both displayed integrals are absolutely convergent, and
$C_{m,n}(x)=(G_m*G_n)(x)$. In particular, every expression
\eqref{stconv:eq:Cmn} has a finite exact reduction using Theorem~\ref{stconv:thm:product}.
\end{theorem}
\begin{proof}
Near $t=0$, the shift relation~\eqref{stconv:eq:stshift} leaves a possible
singularity
\[
 a_m(t)[\gamma_n(x-t)-\gamma_n(x)]=O(|\log t|^m).
\]
The other endpoint is the same argument with $m,n$ interchanged. The
second integral is nonsingular. These observations prove absolute
convergence, locally uniformly with derivatives in $x$ after using
endpoint-fixed local coordinates.

For the normalization, approximate $G_m$ by the cutoff function
$\mathbf1_{[\epsilon,1)}\gamma_m$ plus
$\log^{m+1}\epsilon\,\delta_0/(m+1)$, and similarly for $G_n$. On a
compact subinterval of $0<x<1$, their convolution consists of
\[
 \int_\epsilon^{x-\epsilon}\gamma_m(t)\gamma_n(x-t)\dd t
 +\int_x^1\gamma_m(t)\gamma_n(1+x-t)\dd t
\]
plus the two logarithmic counterterms multiplying $\gamma_n(x)$ and
$\gamma_m(x)$. The product of the two delta counterterms is supported
at $0$, outside that compact subinterval. Integrating the two subtracted
$a$-terms explicitly and taking the limit gives~\eqref{stconv:eq:Cmn}.
The locally uniform endpoint argument identifies this limit with the
restriction of the distributional convolution.
\end{proof}

\begin{corollary}[Four scalar identity families]\label{stconv:cor:scalar}
For $0<x<1$,
\begin{align}
 C_{0,0}(x)&=2\gamma_1(x)+\zeta(2),\label{stconv:eq:C00}\\
 C_{0,1}(x)&=\frac32\gamma_2(x)-\zeta(2)\gamma_0(x)-\zeta(3),\label{stconv:eq:C01}\\
 C_{0,2}(x)&=\frac43\gamma_3(x)-2\zeta(2)\gamma_1(x)
                    +2\zeta(3)\gamma_0(x)+2\zeta(4),\label{stconv:eq:C02}\\
 C_{1,1}(x)&=\gamma_3(x)-2\zeta(2)\gamma_1(x)
                    +2\zeta(3)\gamma_0(x)+\frac14\zeta(4).
 \label{stconv:eq:C11}
\end{align}
\end{corollary}
\begin{proof}
Restrict~\eqref{stconv:eq:product00}--\eqref{stconv:eq:product11} to $0<x<1$, where
$G_n$ agrees with $\gamma_n$ and $\stconvId$ agrees with the constant $-1$.
\end{proof}

\subsection{A concrete digamma integral}
Since $\gamma_0=-\psi$, the first case reads
\begin{align}
 &\int_0^x\left[\psi(t)\psi(x-t)+\psi(x)\left(\frac1t+\frac1{x-t}\right)\right]\dd t
 \label{stconv:eq:digamma-integral}\\
 &\qquad+\int_x^1\psi(t)\psi(1+x-t)\dd t-2\psi(x)\log x
 =2\gamma_1(x)+\frac{\pi^2}{6}.\notag
\end{align}
The two divergent-looking summands inside the first bracket must be
integrated together; their sum is integrable. Equation~\eqref{stconv:eq:digamma-integral}
is not an assertion that either unregularized convolution integral exists.

At $x=1/2$, the Hurwitz identity
$\zeta(s,1/2)=(2^s-1)\zeta(s)$ implies
\[
 \gamma_1(1/2)=\gamma_1-2\gamma\log2-\log^22.
\]
Thus the left side of~\eqref{stconv:eq:digamma-integral} at $1/2$ equals
\begin{equation}\label{stconv:eq:half}
 2\gamma_1-4\gamma\log2-2\log^22+\frac{\pi^2}{6}.
\end{equation}
The same method at any rational argument turns the full multiplication
table into exact relations among rational-argument Stieltjes functions,
with no integer-relation search.
